% Figure: the friction circle. The tyre can supply a horizontal force
% in any direction up to a magnitude mu*N; the accessible forces fill a
% disc. Braking, cornering, and combined manoeuvres sit inside it.
\begin{figure}[htbp]
\centering
\begin{tikzpicture}[
  vec/.style={draw=engblue, -{Stealth}, very thick},
  comp/.style={draw=enggray, thin, dashed},
  lbl/.style={font=\scriptsize},
]
% the friction limit disc, radius 3 = mu*N
\draw[draw=engred, very thick, fill=engred!6] (0,0) circle (3);
% axes: forward (up) is braking/traction, side (right) is cornering
\draw[-{Stealth}, engblack, thick] (0,-3.4) -- (0,3.6)
  node[above, lbl] {longitudinal (brake / drive)};
\draw[-{Stealth}, engblack, thick] (-3.4,0) -- (3.6,0)
  node[right, lbl] {lateral (cornering)};
% a combined-load demand vector: braking + cornering at 40 deg from vertical
\def\ang{40}
\draw[vec] (0,0) -- ({3*sin(\ang)}, {-3*cos(\ang)})
  node[below right, lbl, engblue] {demand $\vec{F}$};
% its components
\draw[comp] ({3*sin(\ang)},{-3*cos(\ang)}) -- (0,{-3*cos(\ang)});
\draw[comp] ({3*sin(\ang)},{-3*cos(\ang)}) -- ({3*sin(\ang)},0);
\node[lbl, enggray, anchor=east] at (0,{-1.5*cos(\ang)}) {$F_{\text{brake}}$};
\node[lbl, enggray, anchor=south] at ({1.5*sin(\ang)},0) {$F_{\text{corner}}$};
% radius label
\draw[-{Stealth}, engamber, thick] (0,0) -- ({3*cos(30)},{3*sin(30)});
\node[lbl, engamber, anchor=south west] at ({1.6*cos(30)},{1.6*sin(30)})
  {$\mu N$};
% a point outside the circle (loss of grip)
\fill[engblack] ({3.4*sin(20)},{3.4*cos(20)}) circle (1.5pt);
\node[lbl, anchor=west] at ({3.4*sin(20)},{3.4*cos(20)})
  {outside: tyre slides};
\end{tikzpicture}
\caption[The friction circle for combined tyre loading]{The friction
circle. A tyre can return a horizontal force in any direction up to a
magnitude $\mu N$, so the forces it can supply fill a disc of radius
$\mu N$ (red). Pure braking spends the whole budget along the vertical
axis; pure cornering spends it along the horizontal axis; a driver who
brakes and turns at once spends it on a slant, and the two demands add
as vectors. A demand vector that reaches the rim uses all the available
grip; one that lies inside has grip to spare; one that would fall
outside cannot be met and the tyre slides. The circle is why trail
braking into a corner must ease the brake as steering builds: the sum
of the two, not either alone, must stay within the rim.}
\label{fig:vol03ch04:friction-circle}
\end{figure}
